% generated by roam.py handbook export -- не редактировать, править заметки
% 20260922174919-связь_сильнои_выпуклости_и_липшицевости_градиента.org (body 119302da06f7)

\begin{theorem}[{Связь сильной выпуклости и липшицевости градиента}]\label[theorem]{note:3C6CF60E-0773-488E-A0B9-4D7510D456C0}
  Пусть $f$~--- $\mu$-\hyperref[note:22E546FC-186F-4A67-85AA-2C85DCEF9021]{сильно выпукла} и $f \in C^{1,1}_L$~--- функция с \hyperref[note:7DD1C7AB-1A49-4675-9F21-5A16C717F670]{липшицевым градиентом}.
  Тогда $\mu \le L$
\end{theorem}

\begin{remark}
  Зачем: для градиентного спуска $\| x_k - x^{*} \|^2 \le \left( 1 - \frac{\mu}{L} \right)^k \| x_0 - x^{*} \|^2$, и $\mu \le L$ гарантирует $1 - \frac{\mu}{L} \in [0;1)$. Отношение $\kappa = \frac{L}{\mu} \ge 1$ называют числом обусловленности: чем оно больше, тем медленнее сходимость.
\end{remark}

\begin{remark}
  Интуиция: сильная выпуклость означает, что функция ограничена параболой снизу, а липшицевость градиента~--- что функция ограничена параболой сверху. Тогда $\mu,L$ отвечают за растянутость этих парабол, следовательно парабола снизу должна быть не уже, чем сверху, отсюда и $\mu \le L$.
\end{remark}

\begin{proof}
  Давайте обозначим $r(x;y) := f(y) - f(x) - \left\langle \nabla f(x); y-x \right\rangle$~--- остаток аппроксимации первого порядка.

  Пусть $E$~--- область определения $f$.
  По \hyperref[note:1B5A3E25-56E8-4D5E-9CF6-E2731DA77604]{дифференциальному критерию сильной выпуклости}
  \[
  \forall x,y \in E\ f(y) \ge f(x) + \left\langle \nabla f(x); y -x \right\rangle + \frac{\mu}{2} \| y -x \|^2
  \]
  Давайте перенесем некоторые слагаемые налево, получим
  \[
  \forall x,y \in E\ f(y) - f(x) - \left\langle \nabla f(x); y -x \right\rangle \ge \frac{\mu}{2} \| y -x \|^2
  \Leftrightarrow r(x;y) \ge \frac{\mu}{2} \|y-x\|^2
  \]

  Теперь оценим $r$ сверху через \hyperref[note:379E4CAD-A44B-4631-8DAF-F3DC708EAEDE]{усиленное неравенство для функций с липшицевым градиентом}
  \[
  \forall x,y \in E\ \left| f(y) -f(x) - \left\langle \nabla f(x); y-x  \right\rangle \right|\le \frac{L}{2} \| y - x\|^2_2
  \]
  Это равносильно
  \[
  \forall x,y \in E\ - \frac{L}{2} \| y - x\|^2_2 \le  f(y) -f(x) - \left\langle \nabla f(x); y-x  \right\rangle \le \frac{L}{2} \| y - x\|^2_2
  \]
  Воспользуемся только правой частью, получим
  \[
  r(x;y) \le \frac{L}{2} \| y-x \|^2
  \]
  Итого
  \[
  \frac{\mu}{2} \| y-x \|^2 \le r(x;y) \le \frac{L}{2} \| y-x\|^2 \Rightarrow \frac{\mu}{2} \| y-x\|^2 \le \frac{L}{2} \|y-x\|^2
  \]
  Давайте возьмем $y \ne x$ и разделим обе части на $\frac{1}{2} \| y-x \|^2$.
  Получим $\mu \le L$.
\end{proof}
